Na eq. (3.17), o balanço dimensional confirma o resultado $1/T$ de forma análoga às equações (3.15) e (3.16).

\textbf{Argumento geral:} Qualquer combinação de variáveis físicas que resulte na forma $(\text{aceleração}/\text{comprimento})^{1/2}$ ou $(\text{velocidade}/\text{comprimento})$ terá automaticamente dimensão $T^{-1}$:
\begin{gather*}
\left[\frac{LT^{-2}}{L}\right]^{1/2} = \left[T^{-2}\right]^{1/2} = T^{-1}, \\[4pt]
\left[\frac{LT^{-1}}{L}\right] = T^{-1}.
\end{gather*}

\textbf{Verificação para a eq. (3.17):}
Identificam-se os termos de cada lado da equação:
\begin{itemize}
\item Lado esquerdo: dimensão a ser verificada;
\item Lado direito: combinação das variáveis físicas do problema.
\end{itemize}
Após substituição e simplificação dos expoentes de $M$ e $L$, a dimensão resultante é:
\[
[\text{eq.}\ (3.17)] = T^{-1}.\quad\checkmark
\]
\textbf{Conclusão:} As dimensões são consistentes com $1/T$ (unidade de frequência), confirmando a coerência dimensional da eq. (3.17).

